\section{Division in the spectral dual cokernel and exact
finite-dimensional
lifting}\label{division-in-the-spectral-dual-cokernel-and-exact-finite-dimensional-lifting}

24 September 2026. Independent derivation SCL0--SCL9.

This proves that every \(L^t-a\) acts bijectively on the actual
algebraic cokernel \(\mathscr C=Q'/\mathcal D_LQ\). It constructs
division on the original Fréchet quotient before taking transposes, then
gives the explicit unique lift of every finite-dimensional
\(L^t\)-invariant subspace of \(Q'\). The full real-group covariance of
those finite lifts is proved from the original jet action. No assertion
that the whole cokernel vanishes, no geometric purity, and no RH
conclusion is inserted.

\subsection{SCL0. Source definitions, prior results and operation
prerequisites}\label{scl0.-source-definitions-prior-results-and-operation-prerequisites}

The controlling source remains B1--B5 and their proved consequences in
\href{../foundations/user_definitions_20260924/SOURCE_OPERATIONS_AND_PROOFS.md}{\nolinkurl{SOURCE_OPERATIONS_AND_PROOFS.md}},
with the verbatim U01--U18 in
{USER\_DEFINITIONS\_VERBATIM.md}.
The complete global user arguments USR-9ad1c0a2d09dba92,
USR-f55d16feb948d8c2 and USR-6152e3bc6302258c, and the operation
instruction USR-4322be19bff532cd, remain controlling as recorded in
\nolinkurl{CORPUS_AND_OPERATION_RULES.md} (private source record) and ASD0. They are not replaced by
an assumption about a finite prime system.

The current calculation uses already reconstructed integer arithmetic
and the original \(\zeta\). \(Z_0\) retains absence, \(Z_1/\tau\)
retains primitive presence without \(Z_2\) parity, and the original
integer layer retains its supplied amount and parity data. There is no
source addition on \(\tau\). A continuous complex dual, a polynomial in
\(L\), and a receiving scalar below belong to the constructed
coefficient category, not to an unstated operation on primitive
\(\tau\).

The actual coefficient complex, its faithful source-ring action, and its
dual cone are \url{ACTUAL_SUPPORTED_DUALITY_INDEPENDENT.md},
ASD1--ASD14. The new full residue pairing is
\url{GLOBAL_ORIGINAL_ZETA_RESIDUE_DUALITY_INDEPENDENT.md}, GZR2--GZR8.
The complete jet isolators and projectors are
\href{https://github.com/KokunoYumeto/zeta-function-research-reader/blob/main/workbenches/splitzero-tandem/continuations/20260924-cohomology-weight-and-character-lifts/gct-weight-control/proofs/ORIGINAL_ZETA_RESOLVENT_AND_PROJECTORS.md}{the
retained \nolinkurl{ORIGINAL_ZETA_RESOLVENT_AND_PROJECTORS.md}}, RZ5--RZ6, read
completely for this derivation; GZR4 and GZR7--GZR8 were also read. The
full Schwartz image theorem and original Mellin topology are
\url{ORIGINAL_MELLIN_SPECTRAL_SYNTHESIS.md}, OMS1--OMS5.

Human provenance remains Alain Connes and Caterina Consani,
\href{https://arxiv.org/abs/0903.2024v3}{\emph{Schemes over
\(\mathbb F_1\) and zeta functions}, arXiv:0903.2024v3, §5}, for the
coefficient sheaf and Fourier localization; Ralf Meyer,
\href{https://arxiv.org/abs/math/0412277v3}{\emph{A spectral
interpretation for the zeros of the Riemann zeta function},
arXiv:math/0412277v3}, for the analytic closed-image framework, with the
programme's exact original-space return proved in the cited sources; and
Pierre Deligne,
\href{https://www.numdam.org/item/PMIHES_1980__52__137_0/}{\emph{La
conjecture de Weil. II}, §3.6}, for the lifting-obstruction argument
motivating this continuation. The new calculation does not attribute a
numerical purity theorem to these receiving maps.

\subsection{SCL1. The exact spaces, auxiliary source transform and
original jet
representatives}\label{scl1.-the-exact-spaces-auxiliary-source-transform-and-original-jet-representatives}

Let \[
 \mathcal B=\{F\text{ entire}:
 b_{A,N}(F)=\sup_{|\sigma|\le A,t\in\mathbb R}
 (1+|t|)^N|F(\sigma+it)|<\infty\quad(A,N\ge0)\},
\] \[
 I_\zeta=\{F\in\mathcal B:F^{(j)}(\rho)=0
 \text{ for every actual nontrivial zero }\rho,\
 0\le j<m_\rho\},\qquad Q=\mathcal B/I_\zeta .
 \tag{SCL1.1}
\] Use the original Fréchet topology and quotient topology. The ideal is
closed by Cauchy's derivative estimates. The original raw test-space map
is \[
 \mathcal M_0b(s)=\int_0^\infty b(u)u^s\,\frac{du}{u},\qquad
 \mathcal M_0\Sigma h(s)
 =2\zeta(s)\int_0^\infty h(v)v^s\,\frac{dv}{v}.
 \tag{SCL1.2}
\] The cited complete image theorem identifies \(Q\) with the actual
\(A/\Sigma S\), not with an unrestricted formal jet product. The earlier
centered receiver retains the exact comparison \[
 \mathcal Tk(u)=2u^{-1/2}k(u),\qquad
 \mathcal T^{-1}b(u)=\tfrac12u^{1/2}b(u),\qquad
 \mathcal M_0\mathcal Tk=2F_k.
 \tag{SCL1.3}
\] No factor of this comparison is changed in the calculation below.

Write \(L[F]=[sF]\), a continuous operator. Let \(Q'\) be its continuous
complex-linear dual, and \[
 (L^t\lambda)(q)=\lambda(Lq).
\] All transposes preserve the actual continuous dual. For exact dual
topology statements use \(\sigma(Q',Q)\); the displayed transposes are
also strong-dual continuous by ASD2.2.

The auxiliary source function, used with its full multiplier, is \[
 F_*(s)=\frac{s(s-1)}8\pi^{-s/2}\Gamma(s/2)\zeta(s)\in\mathcal B.
 \tag{SCL1.4}
\] It is the transform of the explicit original source
\(f_*(v)=\frac{\pi}{2}v^2(2\pi v^2-3)e^{-\pi v^2}\). Its zero divisor is
exactly the actual nontrivial-zero divisor with every multiplicity. Its
retained exceptional values include \[
 F_*(0)=F_*(1)=\frac18,
\] \[
 F_*(-2r)=
 \frac{(1+2r)2r}{8}\pi^{-(1+2r)/2}
 \Gamma((1+2r)/2)\zeta(1+2r)\ne0\qquad(r\ge1).
 \tag{SCL1.5}
\] These are the original multiplier's removable values, proved in OMS2.
They are not discarded by calling \(F_*\) an auxiliary function.

For an actual zero \(a\) of order \(m\), put \[
 A_a(s)=\frac{F_*(s)}{(s-a)^m}.
\] It is entire and in \(\mathcal B\), and \(A_a(a)\ne0\). Recursively
compute the first \(m\) Taylor coefficients of \(1/A_a\) at \(a\),
obtaining a polynomial \(B_a\) such that \[
 e_a(s)=A_a(s)B_a(s)=1+O((s-a)^m).
\] Explicitly, if \(A_a(a+t)=\sum_{r\ge0}a_rt^r\), set \(b_0=a_0^{-1}\)
and \(b_r=-a_0^{-1}\sum_{h=1}^ra_hb_{r-h}\), and take
\(B_a(a+t)=\sum_{r<m}b_rt^r\). Then \[
 E_{a,j}(s)=e_a(s)(s-a)^j,\qquad
 e_{a,j}=[E_{a,j}]\in Q,\qquad
 \ell_{a,j}([F])=\frac{F^{(j)}(a)}{j!}.
 \tag{SCL1.6}
\] For \(0\le j<m\) the function \(E_{a,j}\) has Taylor coefficient1 in
degree \(j\) at \(a\), its other coefficients below \(m\) are zero, and
every required jet at every other zero is zero. This follows from the
reciprocal recursion and the full zero divisor of \(F_*\). In particular
\[
 \ell_{a,i}(e_{a,j})=\delta_{ij},\qquad
 (L-a)e_{a,j}=e_{a,j+1}\ (j<m-1),\qquad
 (L-a)e_{a,m-1}=0.
 \tag{SCL1.7}
\] The continuous projector \[
 P_aq=\sum_{j<m}\ell_{a,j}(q)e_{a,j}
 \tag{SCL1.8}
\] has finite-dimensional image \(Q_a\), gives \(Q=Q_a\oplus\ker P_a\)
topologically, and retains every other primary observation. Equality of
all original zero jets implies equality in \(Q\) by SCL1.1; no infinite
sum of projectors is used.

\subsection{SCL2. Construct division on the original source before
taking a
transpose}\label{scl2.-construct-division-on-the-original-source-before-taking-a-transpose}

For \(a\notin\mathscr Z\), the actual set of nontrivial zeros,
\(F_*(a)\ne0\). Define the entire source operator \[
 \mathscr S_aF(s)=
 \frac{F(s)-F(a)F_*(s)/F_*(a)}{s-a}.
 \tag{SCL2.1}
\] The numerator vanishes at \(a\). This is continuous on
\(\mathcal B\). To verify the division estimate, outside a fixed disk
about \(a\), division by \(s-a\) is bounded on every bounded vertical
strip at bounded height and is \(O((1+|\Im s|)^{-1})\) at large height.
Inside that disk the quotient is holomorphic and bounded by a Cauchy
estimate on a larger fixed disk. Evaluation \(F(a)\) and all derivatives
on these fixed disks are controlled by a finite strip seminorm. These
estimates prove every output seminorm is controlled by finitely many
input seminorms.

If \(F\in I_\zeta\), both numerator terms have all required zero jets.
Since \(a\) is not one of those zeros, division preserves them.
Therefore \(\mathscr S_a\) descends to a continuous \(S_a:Q\to Q\), and
direct multiplication gives \[
 (L-a)S_a=S_a(L-a)=\operatorname{id}_Q.
 \tag{SCL2.2}
\] The term involving \(F_*\) vanishes in the quotient; it remains
present in the entire source formula SCL2.1, including at
\(a=0,1,-2,-4,\ldots\).

For an actual zero \(a\) of order \(m\), the naive division by \(s-a\)
needs its full source correction. Set \[
 H_F(s)=F(s)-F(a)E_{a,0}(s)
\] and define \[
 \boxed{
 \mathscr S_aF(s)=
 \frac{H_F(s)-F_*(s)H_F^{(m)}(a)/F_*^{(m)}(a)}
      {s-a}.}
 \tag{SCL2.3}
\] The denominator \(F_*^{(m)}(a)\) is nonzero by the actual
multiplicity. The numerator vanishes at \(a\), so the same division
estimate proves an entire member of \(\mathcal B\), continuously
depending on \(F\). Every evaluation and derivative in SCL2.3 is a
continuous finite-order evaluation on the original space.

If \(F\in I_\zeta\), then \(F(a)=0\), and subtraction of the displayed
multiple of \(F_*\) cancels its Taylor coefficient of order \(m\) at
\(a\). The resulting numerator vanishes to order at least \(m+1\) there,
and to each required order at every other zero. Division by \(s-a\)
consequently leaves a member of \(I_\zeta\). Thus SCL2.3 descends to a
well-defined continuous \(S_a:Q\to Q\).

This proves why the correction is part of the actual map. Without it,
dividing an element of \(I_\zeta\) of exact order \(m\) at \(a\)
produces a nonzero \((m-1)\)-jet in the quotient.

The two exact inverse identities are \[
 \boxed{
 (L-a)S_a=1-e_{a,0}\ell_{a,0},\qquad
 S_a(L-a)=1-e_{a,m-1}\ell_{a,m-1}.}
 \tag{SCL2.4}
\] For the first, multiply SCL2.3 by \(s-a\); the \(F_*\) term lies in
\(I_\zeta\), leaving \([F]-F(a)e_{a,0}\). For the second, apply SCL2.3
to \((s-a)F\). Its value at \(a\) is zero and its \(m\)-th derivative is
\(mF^{(m-1)}(a)\). Moreover \[
 \left[\frac{F_*(s)}{s-a}\right]
 =\frac{F_*^{(m)}(a)}{m!}\,e_{a,m-1},
\] because these are exactly its full jets at all actual zeros. The
correction therefore subtracts \(F^{(m-1)}(a)e_{a,m-1}/(m-1)!\), proving
the second identity with all factorials.

Also \(S_ae_{a,0}=0\), by inserting the specific source representative
\(E_{a,0}\) in SCL2.3. Together with SCL2.4 this gives \[
 S_ae_{a,j}=e_{a,j-1}\quad(1\le j<m).
 \tag{SCL2.5}
\] On \(\ker P_a\), the same source formula has zero \(a\)-jet after the
correction; hence \(S_a\) preserves that subspace. There \(S_a\) is the
inverse of \(L-a\).

It follows constructively that \[
 \ker(L-a)=\mathbb C e_{a,m-1},\qquad
 (L-a)Q=\ker\ell_{a,0}\quad(a\in\mathscr Z).
 \tag{SCL2.6}
\] The range inclusion follows by evaluation at \(a\), the reverse
inclusion by the first identity in SCL2.4. The kernel inclusion follows
by the second identity, and SCL1.7 proves equality. The range is closed
and of codimension1; its continuous right inverse is \(S_a\) restricted
to that range. Thus \(L-a\) is strict onto its image, rather than only
an abstract surjection.

More generally, for \(k\ge1\) and \(r=\min(k,m)\), \[
 \ker(L-a)^k
 =\operatorname{span}\{e_{a,m-r},\ldots,e_{a,m-1}\},
\] \[
 (L-a)^kQ
 =\bigcap_{j=0}^{r-1}\ker\ell_{a,j}.
 \tag{SCL2.7}
\] On \(Q_a\) these are the full truncated multiplication and
backward-shift calculation from SCL1.7 and SCL2.5. On \(\ker P_a\) both
operators are inverses, as just proved. The direct sum in SCL1.8 proves
SCL2.7 on all of \(Q\), with continuous right inverse \(S_a^k\) on the
displayed closed range. No global spectral expansion is required.

\subsection{SCL3. Exact continuous transpose kernels and
ranges}\label{scl3.-exact-continuous-transpose-kernels-and-ranges}

Let \(A_a=L-a\). For \(a\notin\mathscr Z\), transposition of SCL2.2
gives an everywhere-defined continuous inverse \(S_a'\) of \(L^t-a\) on
\(Q'\). Its weak-* and strong continuity follow by transposing a
continuous map on \(Q\).

For \(a\in\mathscr Z\), put \(k_a=e_{a,m-1}\). The exact result is \[
 \boxed{
 \ker(L^t-a)=\mathbb C\ell_{a,0},\qquad
 (L^t-a)Q'=\{\lambda\in Q':\lambda(k_a)=0\}.}
 \tag{SCL3.1}
\] For the kernel, a functional in \(\ker(L^t-a)\) annihilates the range
in SCL2.6, so it is a scalar multiple of the continuous coordinate
\(\ell_{a,0}\). Conversely that coordinate annihilates the range.

The range inclusion follows by evaluating \((L^t-a)\mu\) at
\(k_a\in\ker(L-a)\). For the reverse inclusion, take a continuous
\(\lambda\) with \(\lambda(k_a)=0\). The explicit continuous functional
\(\mu=S_a'\lambda\) satisfies \[
 (L^t-a)S_a'\lambda
 =(S_a(L-a))'\lambda
 =\lambda-\lambda(k_a)\ell_{a,m-1}
 =\lambda.
 \tag{SCL3.2}
\] This constructs the functional in the correct continuous dual,
without replacing it by an algebraic-dual extension.

The range in SCL3.1 is closed in the weak-* topology and in the strong
topology, since evaluation at the fixed vector \(k_a\) is continuous in
both. The right inverse \(S_a'\) is continuous in both topologies. The
transpose is therefore strict onto its image. The construction proves
the claimed prerequisites before using its kernel or cokernel.

The same argument with SCL2.7 gives \[
 \ker(L^t-a)^k
 =\operatorname{span}\{\ell_{a,0},\ldots,\ell_{a,r-1}\},
 \qquad r=\min(k,m).
 \tag{SCL3.3}
\] Indeed it is the annihilator of the displayed closed
finite-codimensional range of \((L-a)^k\). A functional annihilating
that range factors through its finite tuple of jet coordinates, and is
exactly a linear combination of those coordinates. The transpose range
is the annihilator of the displayed primal kernel, with continuous right
inverse \((S_a^k)'\) there. These formulas retain the complete nilpotent
multiplicity.

\subsection{SCL4. The full residue pairing supplies the exact two
missing
coefficients}\label{scl4.-the-full-residue-pairing-supplies-the-exact-two-missing-coefficients}

Let \[
 B_\zeta([F],[G])=\frac1{2\pi i}
 \left(\int_{2-i\infty}^{2+i\infty}
       -\int_{-1-i\infty}^{-1+i\infty}\right)
 \frac{F(s)G(1-s)}{\zeta(s)}\,ds,
\qquad
 \mathcal D_Lx(y)=B_\zeta(x,y).
 \tag{SCL4.1}
\] These are the already proved convergent, continuous and nondegenerate
operations GZR2--GZR8 on the original quotient. In particular
\(\mathcal D_L:Q\to Q'\) is injective and strong-dual continuous, and \[
 L^t\mathcal D_L=\mathcal D_L(1-L).
 \tag{SCL4.2}
\]

For an actual zero \(a\), put \(b=1-a\). The original functional
equation preserves the actual multiplicity \(m=m_a=m_b\). Use the full
local original-zeta unit \[
 \zeta(b+t)=t^m u_b(t),\qquad
 u_b(t)=\sum_{j\ge0}\frac{\zeta^{(m+j)}(b)}{(m+j)!}t^j,\qquad
 c_b=u_b(0)^{-1}=\frac{m!}{\zeta^{(m)}(b)}.
 \tag{SCL4.3}
\] Then the two exact global identities are \[
 \boxed{
 (\mathcal D_Lx)(k_a)
   =(-1)^{m-1}c_b\,\ell_{b,0}(x),\qquad
 \mathcal D_L(k_b)=c_b\,\ell_{a,0}.}
 \tag{SCL4.4}
\] To prove the first, choose the globally isolated representative of
\(k_a\) in the second slot. The only possible nonremovable pole of the
original integrand is at \(s=b\). At that point its reflected jet is
\((-t)^{m-1}\). The residue is therefore \((-1)^{m-1}u_b(0)^{-1}F(b)\).
GZR5's finite-pole contour theorem justifies this residue computation on
the full original contour for every \(F\in\mathcal B\), not just a
finite polynomial. For the second identity choose the isolator of
\(k_b\) in the first slot. Its jet is \(t^{m-1}\), leaving residue
\(u_b(0)^{-1}G(a)\). This proves both formulas, including their
different signs and the original derivative denominator.

In particular the entire transpose kernel in SCL3.1 belongs to
\(\mathcal D_LQ\). Also, the obstruction to solving
\((L^t-a)\lambda=\mathcal D_Lx\) is exactly the single value
\(\ell_{b,0}(x)\), and SCL2.6 identifies its vanishing with an already
constructed primal division by \(L-b\).

\subsection{SCL5. Every linear polynomial is invertible on the actual
algebraic
cokernel}\label{scl5.-every-linear-polynomial-is-invertible-on-the-actual-algebraic-cokernel}

Define \[
 M=\mathcal D_LQ\subset Q',\qquad
 \mathscr C=Q'/M.
 \tag{SCL5.1}
\] This is the actual algebraic cokernel in ASD14. The subspace is
invariant under \(L^t\) by SCL4.2. No closedness or surjectivity of
\(\mathcal D_L\) is assumed. Let \(\mathsf T\) be the operator induced
by \(L^t\) on \(\mathscr C\).

For every \(a\in\mathbb C\), \[
 \boxed{\mathsf T-a:\mathscr C\longrightarrow\mathscr C
       \text{ is bijective}.}
 \tag{SCL5.2}
\]

For surjectivity when \(a\notin\mathscr Z\), every \(\lambda\in Q'\) is
\((L^t-a)S_a'\lambda\), by SCL3.

For surjectivity when \(a\in\mathscr Z\), retain \(b=1-a\), \(m=m_a\)
and \(\kappa_a=(-1)^{m-1}c_b\ne0\). Given \(\lambda\in Q'\), set \[
 x_0=\frac{\lambda(k_a)}{\kappa_a}\,e_{b,0},\qquad
 \lambda_0=\lambda-\mathcal D_Lx_0,\qquad
 \mu=S_a'\lambda_0.
 \tag{SCL5.3}
\] By SCL4.4, \(\lambda_0(k_a)=0\). Equation SCL3.2 then gives
\((L^t-a)\mu=\lambda_0\). Thus \((\mathsf T-a)[\mu]=[\lambda]\), proving
surjectivity with an explicit continuous-dual representative.

For injectivity, suppose \[
 (L^t-a)\lambda=\mathcal D_Lx.
 \tag{SCL5.4}
\] If \(a\notin\mathscr Z\), then \(b=1-a\notin\mathscr Z\) too, by the
original zero symmetry. Set \(y=-S_bx\). It satisfies \((1-a-L)y=x\).
SCL4.2 shows \((L^t-a)(\lambda-\mathcal D_Ly)=0\). The transpose kernel
is zero in this case, so \(\lambda=\mathcal D_Ly\).

If \(a\in\mathscr Z\), evaluate SCL5.4 at \(k_a\). Its left side is
zero. SCL4.4 therefore gives \(\ell_{b,0}(x)=0\), and SCL2.4 makes
\(y=-S_bx\) solve \((1-a-L)y=x\) again. The same subtraction now lies in
\(\ker(L^t-a)=\mathbb C\ell_{a,0}\). SCL4.4 puts that kernel in \(M\)
explicitly: \(\ell_{a,0}=c_b^{-1}\mathcal D_L(k_b)\). Hence
\(\lambda\in M\), proving injectivity.

This also proves that the explicit inverse from SCL5.3 is independent of
the representative of the class \([\lambda]\): it solves a bijective
operator equation in \(\mathscr C\). Nothing in this proof assumes that
\(\mathscr C\) itself is zero.

\subsection{SCL6. The exact remaining module has no finite-dimensional
spectral
part}\label{scl6.-the-exact-remaining-module-has-no-finite-dimensional-spectral-part}

Every nonzero polynomial \(P\in\mathbb C[t]\) factors into a nonzero
constant and linear factors. By SCL5.2 each factor \(P(\mathsf T)\) is
bijective. The operators commute, and their inverses commute too:
multiplying \(AB=BA\) by either inverse verifies this directly.
Consequently \[
 \frac{P(t)}{R(t)}\cdot c
   =P(\mathsf T)R(\mathsf T)^{-1}c,\qquad R\ne0,
 \tag{SCL6.1}
\] defines a \(\mathbb C(t)\)-vector-space structure on \(\mathscr C\).
If two fractions agree, cross multiplication and the common invertible
denominator prove their operators agree. Distributivity and
multiplication follow from commuting polynomial operators. The zero
vector space is allowed; the calculation has not decided its dimension
over \(\mathbb C(t)\).

There is no nonzero finite-dimensional \(\mathsf T\)-invariant subspace
of \(\mathscr C\). Indeed its finite-dimensional characteristic
polynomial annihilates it by the matrix adjugate identity, whereas that
polynomial in \(\mathsf T\) is injective on all of \(\mathscr C\).

There is also no nonzero finite-dimensional \(\mathsf T\)-equivariant
quotient of \(\mathscr C\). If \(U\) is an invariant subspace with
finite-dimensional quotient, its quotient characteristic polynomial
\(P\) gives \(P(\mathsf T)\mathscr C\subset U\). But \(P(\mathsf T)\) is
onto, so \(U=\mathscr C\).

In particular any equivariant linear map from a finite-dimensional
operator space to \(\mathscr C\), or from \(\mathscr C\) to a
finite-dimensional operator space, is zero: its image is such a subspace
or quotient.

The actual degree-one action in ASD14 includes \(\chi\). Its
infinitesimal receiving operator is \(1-\mathsf T\), not \(\mathsf T\).
Every nonzero polynomial in \(1-\mathsf T\) is again a nonzero
polynomial in \(\mathsf T\). All conclusions in this section therefore
apply to that exact twisted action as well. This identifies a property
of the remaining whole object rather than deleting its twist.

\subsection{SCL7. Explicit unique lift of every finite invariant dual
subspace}\label{scl7.-explicit-unique-lift-of-every-finite-invariant-dual-subspace}

Let \(V\subset Q'\) be any finite-dimensional \(L^t\)-invariant
subspace. Its image in \(\mathscr C\) is finite-dimensional and
invariant, so SCL6 proves it is zero. Thus \[
 V\subset\mathcal D_LQ.
\] Since \(\mathcal D_L\) is injective, there is exactly one linear lift
\[
 s_V:V\to Q,\qquad \mathcal D_Ls_V=\operatorname{id}_V.
 \tag{SCL7.1}
\] We now construct that lift explicitly rather than leave it as an
inverse on an unspecified image.

A finite-dimensional \(L^t\)-module decomposes into its generalized
eigenspaces by the coprime factors of its characteristic polynomial. The
decomposition follows from polynomial Bezout identities, which give the
finite primary projectors. SCL3.3 shows that every occurring eigenvalue
\(a\) is an actual nontrivial zero and that its full dual generalized
eigenspace is
\(\operatorname{span}\{\ell_{a,0},\ldots,\ell_{a,m_a-1}\}\). Thus there
is a finite set \(P\subset\mathscr Z\) containing the complete support
of \(V\), and every \(\lambda\in V\) has the unique expansion \[
 \lambda=\sum_{a\in P}\sum_{j=0}^{m_a-1}
            \lambda(e_{a,j})\ell_{a,j}.
 \tag{SCL7.2}
\] The coefficient formula follows from SCL1.7. It is valid on the
entire \(Q\), because the generalized eigenfunctional classification in
SCL3.3 is an equality of continuous functionals, not only of
restrictions to a block.

For each \(a\in P\), put \(b=1-a\), \(m=m_a\), and define the full
receiving jet \[
 P_{b,\lambda}(t)=
 \left[
 u_b(t)\sum_{j=0}^{m-1}
   (-1)^j\lambda(e_{a,j})t^{m-1-j}
 \right]_{\deg<m}
 =\sum_{i=0}^{m-1}p_{b,i}(\lambda)t^i.
 \tag{SCL7.3}
\] The brackets mean the image in the explicitly stated truncated
algebra \(\mathbb C[t]/(t^m)\); every contributing coefficient of
\(u_b(t)=\sum_{r\ge0}\zeta^{(m+r)}(b)t^r/(m+r)!\) is retained. Set \[
 \boxed{
 s_V(\lambda)=\sum_{a\in P}\sum_{i=0}^{m_a-1}
                   p_{1-a,i}(\lambda)e_{1-a,i}.}
 \tag{SCL7.4}
\] It is a finite combination of the actual global representatives from
SCL1.6, hence belongs to the original quotient.

To prove its lift identity for an arbitrary \(y=[G]\in Q\), the
first-slot isolators reduce each term of the global pairing to its
single pole \(b=1-a\), by GZR5. Its residue is \[
 [t^{m-1}]
 \frac{P_{b,\lambda}(t)}{u_b(t)}
       \sum_{j=0}^{m-1}(-1)^j\ell_{a,j}(y)t^j
 =\sum_{j=0}^{m-1}\lambda(e_{a,j})\ell_{a,j}(y).
\] For the last equality, SCL7.3 says that the coefficient in degree
\(m-1-j\) of \(P_{b,\lambda}/u_b\) is \((-1)^j\lambda(e_{a,j})\); the
reflected jet contributes the second \((-1)^j\). Summing over \(P\)
yields exactly \(\lambda(y)\) by SCL7.2. This proves SCL7.1 for the
explicit formula.

The lift is continuous: it has finitely many continuous evaluations of
\(\lambda\) at fixed vectors, followed by a finite linear combination of
fixed vectors of \(Q\). It is therefore continuous for the subspace
weak-* or strong topology on \(V\), and for its usual finite-dimensional
topology. No continuous inverse of \(\mathcal D_L\) on its entire image
is inferred.

Its generator identity is \[
 (1-L)s_V=s_V L^t,\qquad
 Ls_V=s_V(1-L^t).
 \tag{SCL7.5}
\] Apply \(\mathcal D_L\) to the first difference, use SCL4.2 and
SCL7.1, and then use injectivity. The lifted finite space therefore has
the reflected eigenvalues \(1-a\) and the exact negative transpose
nilpotent, not a dropped nilpotent or a parity assigned to \(\tau\).

\subsection{SCL8. Full real and prime covariance proved on the actual
finite
jets}\label{scl8.-full-real-and-prime-covariance-proved-on-the-actual-finite-jets}

For \(r>0\), use the original \(T_r[F]=[r^sF(s)]\) and the
contragredient \(T_r^\vee\lambda=\lambda\circ T_{r^{-1}}\). Their exact
finite jet actions are \[
 T_r^\vee\ell_{a,j}
 =r^{-a}\sum_{k=0}^{j}
       \frac{(-\log r)^{j-k}}{(j-k)!}\ell_{a,k},
 \tag{SCL8.1}
\] \[
 T_re_{b,i}
 =r^b\sum_{k=0}^{m_b-1-i}
       \frac{(\log r)^k}{k!}e_{b,i+k}.
 \tag{SCL8.2}
\] Equation SCL8.1 follows by differentiating \(r^{-s}F(s)\) at \(a\),
with its full factorials. For SCL8.2 multiply the original Taylor series
\(t^i\) by \(r^{b+t}\) at \(b\); all other-zero jets remain zero. The
full zero-jet separation SCL1.1 proves equality in the original \(Q\).
Thus both formulas concern the actual globally represented blocks.

On any finite sum of these blocks, these operators equal, respectively,
\(\exp(-(\log r)L^t)\) and \(\exp((\log r)L)\), where these are
finite-dimensional matrix functions with the nilpotent sums just
displayed. A polynomial in the finite matrix gives each exponential:
prescribe its Taylor coefficients through the relevant Jordan orders at
the finitely many distinct eigenvalues and use finite polynomial
interpolation. Consequently every \(L^t\)-invariant finite subspace
\(V\) above is invariant under \(T_r^\vee\), and its lift \(s_VV\) is
invariant under \(T_r\).

Now SCL7.5 and the actual finite matrix formulas prove \[
 \boxed{T_rs_V=r\,s_VT_r^\vee\quad(r>0).}
 \tag{SCL8.3}
\] Indeed exponentiating the finite identity \(Ls_V=s_V(1-L^t)\) gives
\(T_rs_V=s_V\exp((\log r)(1-L^t))
=r\,s_V\exp(-(\log r)L^t)\), and the preceding direct jet calculation
identifies each exponential with the original group action. This is a
proved whole-group relation on \(V\); no passage from an infinitesimal
relation to an unproved infinite-dimensional exponential has occurred.
In particular it holds at every original prime \(r=p\).

Thus \(s_V:\chi\otimes V\to Q\) is equivariant for exactly the action in
ASD14. The source receiving-ring action on this spectral part factors
through its first scalar coordinate \(\epsilon\), so the complex-linear
map \(s_V\) is also source-equivariant there: every integer \(n\) acts
by \(n\), and \(\tau\) by identity. The extra \(\epsilon=0\) copies
distinguishing \(\tau\) from integer \(1\) remain in the full cone and
are not discarded by this statement about its spectral summand.

Although every \(\mathcal D_LQ\) is invariant under the full
contragredient real group, no unrestricted infinite-dimensional
exponential formula is claimed on \(\mathscr C\). Its
\(\mathbb C(t)\)-module result and the finite-space covariance have
their independently proved domains.

Conversely, a finite-dimensional subspace \(V\subset Q'\) invariant
under the entire receiving real group is \(L^t\)-invariant. To prove the
differentiation prerequisite, for real \(h\) the original transform
satisfies \[
 \frac{e^{-hs}F(s)-F(s)}h\longrightarrow-sF(s)
 \quad\text{in }\mathcal B.
 \tag{SCL8.4}
\] The integral remainder formula for the scalar exponential bounds the
error on \(|\Re s|\le A\) by \(\frac12|h|\,|s|^2e^{A|h|}|F(s)|\). This
is controlled in each \(b_{A,N}\) by a constant times
\(|h|e^{A|h|}b_{A,N+2}(F)\), proving the asserted convergence. It
descends continuously to \(Q\). Thus, for \(\lambda\in V\) and
\(q\in Q\), \[
 \left.\frac{d}{dh}(T_{e^h}^\vee\lambda)(q)\right|_{h=0}
 =-\lambda(Lq)=-(L^t\lambda)(q).
 \tag{SCL8.5}
\] This is differentiation in the weak-* topology of the actual
continuous dual.

The subspace \(V\) is weak-* closed and has its ordinary
finite-dimensional topology. One explicit verification is to take a
basis of \(V\) and finitely many vectors \(q_i\in Q\) whose evaluation
matrix has full rank; such vectors exist because these continuous
functionals are linearly independent. Solving that finite matrix defines
a weak-* continuous projection \(Q'\to V\). Its fixed space is exactly
\(V\), so \(V\) is closed; its finitely many evaluation coordinates give
its ordinary topology. Difference quotients of the orbit in SCL8.5 lie
in \(V\), and closedness puts their limit \(-L^t\lambda\) in \(V\). This
proves the required invariant-subspace assertion. Invariance under
\(rT_r^\vee\) is equivalent to invariance under \(T_r^\vee\), since the
scalar \(r\) is nonzero, so the statement includes the exact
\(\chi\)-twisted receiving group.

Its lift is the same explicit \(s_V\). There is also a direct
whole-group covariance proof using the injection, independent of
differentiating an inverse: by SCL4.2's group version, \[
 \mathcal D_L(T_rs_V\lambda)
 =rT_r^\vee\lambda
 =\mathcal D_L(s_V(rT_r^\vee\lambda)).
\] Injectivity gives exactly SCL8.3. Nothing in this argument says that
the entire \(Q\) or \(Q'\) is finite-dimensional or pure.

\subsubsection{Every receiving support label in these
maps}\label{every-receiving-support-label-in-these-maps}

Retain the receiving lattice \(L_{\rm supp}\) and the carrier proved in
GZR9: \[
 G_{L_{\rm supp}}(W)
 =\{(0,\alpha):\alpha\in L_{\rm supp}\}
   \cup\{(w,1_{L_{\rm supp}}):w\in W\}.
\] For every constructed receiving linear map \(f\) in this note its
exact lift is \[
 G_{L_{\rm supp}}(f)(w,\alpha)=(f(w),\alpha).
 \tag{SCL8.6}
\] It is well-defined: a nontop-labelled input has amplitude zero, and
linearity sends that amplitude to zero; a nonzero output can therefore
occur only at top support. Substitution proves preservation of identity
and composition. In particular \[
 G_{L_{\rm supp}}(\mathcal D_L)
 G_{L_{\rm supp}}(s_V)=\operatorname{id}_{G_{L_{\rm supp}}(V)}.
 \tag{SCL8.7}
\] These formulas apply to the exact original quotient, projectors,
division operators, continuous transpose maps and the finite lifts, as
well as the supported maps explicitly computed in ASD6. A vanishing map
amplitude retains the output \((0,\alpha)\), not an erased support
label. The receiving zero \((0,\alpha)\) is not identified with
primitive \(\tau\).

Independently labelled blocks remain separate factors of the
corresponding products of these carriers, with the block lift applied to
its own factor. No identification of \(G_{L_{\rm supp}}(W_1\oplus W_2)\)
with \(G_{L_{\rm supp}}(W_1)\times G_{L_{\rm supp}}(W_2)\) is made: the
former has a common label, whereas the latter retains two labels. This
spells out the actual receiving lifts without assigning source \(\tau\)
a scalar amplitude or asserting an abelian category of primitive
\(\tau\) objects.

\subsection{SCL9. Exact consequence for the actual supported
dual-comparison
cone}\label{scl9.-exact-consequence-for-the-actual-supported-dual-comparison-cone}

The actual ASD14 chain map is \[
 \Psi_\zeta:D_{\rm full}\to\chi\otimes D_{\rm full}^\vee[-2],
 \qquad \Psi_\zeta^1=\pi'\mathcal D_L\pi,\quad\Psi_\zeta^0=0.
\] Its cohomology sequence contains the exact global row \[
 \boxed{
 0\longrightarrow Q\xrightarrow{\mathcal D_L}
 \chi\otimes Q'\xrightarrow{q_{\rm cone}}
 \chi\otimes\mathscr C\longrightarrow0.}
 \tag{SCL9.1}
\] The target actions retain the twist. Under the exact identification
of ASD14.7, \(q_{\rm cone}\) is the quotient map
\(\lambda\mapsto[\lambda]\), induced by the comparison triangle's map
from its target to its cone. It is not the connecting arrow. The actual
following connecting arrow has domain \(\chi\otimes\mathscr C\) and
codomain \(H^2(D_{\rm full})=0\), so that following arrow is zero.

For every finite-dimensional \(L^t\)-invariant subspace \(V\subset Q'\),
the map to cone cohomology is zero on \(\chi\otimes V\), and the exact
lift through the first arrow of SCL9.1 is uniquely the map \(s_V\) in
SCL7.4. This is proved by SCL7.1, and its full real/prime equivariance
is SCL8.3. Thus the actual comparison lifts every finite spectral
subspace, with all Jordan orders retained. This conclusion improves the
earlier statement that only a weak-* dense image was known.

The remaining degree-one cone group is a module on which every nonzero
polynomial in its actual generator \(1-L^t\) is invertible. It has
neither a nonzero finite-dimensional invariant subspace nor a nonzero
finite-dimensional equivariant quotient. It is not proved to vanish as
an algebraic group; its weak-* Hausdorff quotient was already proved
zero in ASD14.8, and that topological statement is not substituted for
the new algebraic calculation.

All other cone terms remain exactly as computed: \[
 H^{-1}(K_\zeta)=H_0\oplus V_{\rm extra},\qquad
 H^0(K_\zeta)=0,\qquad
 H^2(K_\zeta)=\chi\otimes(H_0'\oplus V_{\rm extra}').
 \tag{SCL9.2}
\] In particular the four original endpoint lines, their duals, and both
full faithful closed coefficient copies remain. The unique lift
constructed here lands in \(Q=H^1(D_{\rm full})\). It is not an asserted
equivariant cochain section \(Q\to A\); the finite source
representatives in SCL7.4 are supplied explicitly without that stronger
claim.

This is a concrete lifting theorem for the actual supported
dual-comparison obstruction. It neither identifies that comparison with
Deligne's full geometric specialization nor supplies Deligne's required
numerical upper and lower weight bounds. The result says exactly what
the surviving obstruction can no longer contain, and constructs the
unique lifts for the finite spectral objects it excludes.
